\section{Where the Registration Comes From}
\label{sec:ch02-where-the-registration-comes-from}

\index{source generator|(}%
You can read the method you just called. Build with the generated sources kept
on disk, into a directory outside the project so that the next build does not
compile them a second time.

\begin{minted}[fontsize=\footnotesize, breakanywhere]{text}
dotnet build -p:EmitCompilerGeneratedFiles=true -p:CompilerGeneratedFilesOutputPath=../gen
\end{minted}

Two generators write into that path, and one of them is the ASP.NET Core SDK's
own. Hot Chocolate's sits under a directory named for the analyzers package.
Below is its body, with the signature broken across lines and the
\code{global::} qualifications on Hot Chocolate's own types dropped, both to
fit the page; above it the real file carries a using block and a pair of
pragmas.

\begin{minted}[fontsize=\footnotesize]{csharp}
public static partial class SessionsTypesRequestExecutorBuilderExtensions
{
    public static IRequestExecutorBuilder AddSessionsTypes(
        this IRequestExecutorBuilder builder)
    {
        builder.AddTypeExtension(typeof(global::Sessions.Query));
        builder.ConfigureSchema(
            b => b.TryAddRootType(
                () => new ObjectType(
                    d => d.Name(OperationTypeNames.Query)),
                OperationType.Query));
        return builder;
    }
}
\end{minted}

\index{schema!built at compile time}%
The registration list is a compile-time artifact. \code{[QueryType]} is not
read at startup by a reflection scan over your assembly. It is a signal to a
Roslyn generator, which emits a method naming each type explicitly. So the
attribute is not documentation. It is the whole of the wiring, and the
consequence of leaving it off depends entirely on whether it was the last one.

Take the attribute off a second class that was going to add fields, while some
other class still carries one, and the project compiles, the service starts,
the requests you already had keep working, and the fields that class would have
contributed are absent from the schema. Nothing is logged. Take it off the only
class that has one and you get the opposite, a startup failure that tells you
exactly what is wrong:

\begin{minted}[fontsize=\footnotesize]{text}
HotChocolate.SchemaException: The schema builder was unable to identify
the query type of the schema. Either specify which type is the query type
or set the schema builder to non-strict validation mode.
\end{minted}

That one stops the host before it listens, so you meet it the moment you press
run. The silent one you meet when a client asks for a field that is not there.

The method's name comes from the assembly name and nowhere else. Rename the
project and this method is renamed with it, and the call in \code{Program.cs}
stops compiling. That is a reasonable trade, since a broken build beats an
empty schema, but it does mean that the registration line you copy from any
tutorial is wrong for your project by construction.

\index{Query type}%
The root \code{Query} type is created by the generator rather than by you. Your
class is registered as a \emph{type extension} on it. That is how a static
class with no base type and no interface contributes fields to a type it never
mentions, and it is also how several classes contribute to the same
\code{Query} without any of them owning it. Hold on to that shape. When
chapter~\ref{ch:first-subgraph} turns this service into a subgraph, the
federation library adds its own fields to that same \code{Query} by the same
mechanism, and \code{\_service} and \code{\_entities} will arrive from nowhere
in exactly the way \code{sessions} does now.
\index{source generator|)}%
